\section{Models, shared decoders, and prior-spine benchmarks}\label{sec:model}

\subsection{Matrix states, cyclic shifts, and bounded-arity shared set codes}\label{sec:defs-shared}

For $d \ge 1$, write $\Z_d = \Z/d\Z$ (formalized as \leanname{ZMod d}). Given an entry alphabet
$\Sigma$, a $d \times d$ matrix state is a map $X : \Z_d \to \Z_d \to \Sigma$, written
$X \in \Sigma^{d\times d}$ (\leanname{SharedDecoders.State}). Its transpose $\tau(X) = X^\top$
(\leanname{SharedDecoders.swap}) and column-cyclic shift $\rot_c(X)$ for $c \in \Z_d$
(\leanname{SharedDecoders.rot}) are defined by
\[
  \tau(X)_{i,j} = X_{j,i},
  \qquad
  \rot_c(X)_{i,j} = X_{i,\,j+c}
  \qquad (i, j \in \Z_d).
\]
Thus column $0$ of $\rot_c(X)$ is column $c$ of $X$, whereas column $c$ of $\tau(X)$ is the
transposed vector $r \mapsto X_{c,r}$ (row $c$ of $X$).

\begin{definition}[Syntactic read footprints and bounded-arity shared set codes]\label{def:bounded-arity-set-code}
Let $d \ge 1$, let $M \ge 0$ be a number of decoder classes, let $\Delta \ge 0$ be a read-arity
bound, and let $\Sigma$ be a finite alphabet.
\begin{enumerate}[label=(\alph*)]
\item A family of class decoders $f : \{0,\dots,M-1\} \to \Sigma^{d\times d} \to \Sigma^{\Z_d}$ has
  \emph{syntactic cell support $\le \Delta$} with footprint table
  $Q : \{0,\dots,M-1\} \to \Z_d \to \binom{\Z_d \times \Z_d}{\le \Delta}$
  (\lean{SyntacticCellSupport}) if $\abs{Q_{m,r}} \le \Delta$ for every class $m$ and output
  coordinate $r \in \Z_d$, and whenever $Y_1, Y_2 \in \Sigma^{d\times d}$ agree on every cell
  $q \in Q_{m,r}$, one has $f_m(Y_1)_r = f_m(Y_2)_r$. (Its whole-column counterpart reading at
  most $\Delta_c$ columns of $Y$ is \lean{SyntacticColumnSupport}.)
\item Given a column-class assignment $p : \Z_d \to \{0,\dots,M-1\}$, a broken-output set
  $\mathcal{E} \subseteq \Z_d \times \Z_d$, and a subset $C \subseteq \Sigma^{d\times d}$, the data
  $(\Delta, p, \mathcal{E}, f, Q)$ form a \emph{bounded-arity shared set code on $C$}
  (\lean{IsBoundedAritySetCode}) if $(f, Q)$ has syntactic cell support $\le \Delta$ and
  \[
    f_{p(c)}\bigl(\rot_c(X)\bigr)_r \;=\; X_{c,r}
    \qquad
    \text{for all } X \in C \text{ and all } (r,c) \notin \mathcal{E}.
  \]
\end{enumerate}
\end{definition}

When $\Sigma = V$ is an $F$-module (typically $V = F^b$ for $b$-bit or $b$-symbol lanes),
\leanname{SharedDecoders.IsSharedCode} (\cite[Definition~2.1]{SharedDecoders2026}) restricts the
source $C \le \Mat_{d\times d}(V)$ to an $F$-linear submodule while leaving the class maps $f_m$
arbitrary, whereas linear class-uniform rules (\leanname{SwapRigidity.IsClassUniformCode} and
\leanname{SharedDecoders.LinearClassUniform}, \cite[Definition~6.3]{SeventeenthPaper2026}) further
require each $f_m$ to be $F$-linear, without a cell-arity cap unless explicitly imposed.

\subsection{Native multi-tape Turing machines, canonical words, and frozen cuts}\label{sec:defs-tm}

We use the frozen multi-tape Turing machine model \leanname{Transpose.TM} imported unchanged from
the foundational Lean package of \cite[Definition~2.1]{FirstPaper2026}: a machine $M$ has state
set $\mathrm{Fin}(M.q + 1)$, work-tape alphabet $\mathrm{Fin}(M.g + 1)$, $M.k$ read-write work
tapes, one two-way read-only binary input tape of length $L = d^2 b$ where
$b = \leanname{Transpose.entryBits}(d)$, and one one-way write-only binary output tape.
The canonical word alphabet (\lean{CanonicalWordAlphabet}) is
$\Sigma_d := \{0,1\}^b$ (\leanname{Fin (Transpose.entryBits d) -> Bool}), of exact cardinality
$\abs{\Sigma_d} = 2^b \ge d$.

\begin{lemma}[Canonical word width and dummy-tape padding]\label{lem:pad-dummy-tape}
\formal{NativeSweepInterfaceProofs.lean}{entryBits_eq_clog}
\formal{MachineSweepConsequences.lean}{padDummyWorkTape_transposesWithin}
For every $d \ge 2$, $\leanname{Transpose.entryBits}(d) = \clog{d}$ and
$d \le \abs{\Sigma_d} = 2^{\clog{d}} < 2d$. Moreover, for every time bound $T : \N \to \N$ and
every machine $M$ satisfying \leanname{Transpose.TransposesWithin M T}, there exists a machine
$M^+$ (\lean{padDummyWorkTape_transposesWithin}) with $1 \le M^+.k$ that also satisfies
\leanname{Transpose.TransposesWithin} $M^+$ $T$.
\end{lemma}
\begin{outline}
The identity $\leanname{Transpose.entryBits}(d) = \clog{d}$ is proved from
$\leanname{Nat.clog}$ on binary logarithms, and $\abs{\Sigma_d} = 2^b$ follows by cardinality of
function types over $\mathrm{Fin}\,b$. For dummy-tape padding
(\lean{padDummyWorkTape_transposesWithin}), when $M.k \ge 1$ one sets $M^+ = M$; when $M.k = 0$,
$M^+$ sets $k = 1$, $g = 0$, and $q = M.q$ with one stationary blank work tape whose transition
ignores its read symbol, writes $0$, and stays stationary, maintaining identical control states,
input-head positions, and emitted output bits at every step.
\end{outline}

To interface \leanname{Transpose.TM} both with $F$-linear matrix operators and with nonlinear cut
fibres, we record two exact bridges:
\begin{enumerate}[label=(\alph*)]
\item \emph{Frozen trace/cut fibres} (\lean{IsFrozenTraceCutFibre}): given $M$, dimension $d$,
  horizon $T$, head trajectories $\mathrm{inPos} : \N \to \Z$ and
  $\mathrm{workPos} : \N \to \mathrm{Fin}(M.k) \to \Z$, output emission times
  $\mathrm{emitStep} : \mathrm{Fin}(L) \to \{0,\dots,T-1\}$, and a set of cut steps
  $\Gamma \subseteq \{0,\dots,T-1\}$, a subset $C \subseteq \Sigma_d^{d\times d}$ is a
  \emph{frozen trace/cut fibre} if every $X \in C$ induces the head trajectories
  $(\mathrm{inPos}, \mathrm{workPos})$, emits output cell $(r,c)$'s $\ell$-th bit $X_{c,r,\ell}$ at
  step $\mathrm{emitStep}(\lean{outputBitIdx}(d,r,c,\ell))$, and all $X, X' \in C$ have identical
  post-step control state and written work-tape symbol at the accessed cell of step $u$ for every
  cut step $u \in \Gamma$.
\item \emph{Matrix-to-grid linear bridge}: for every semiring $F$ and $d \ge 0$,
  \lean{encodeTMMatrixToGrid} embeds each Boolean bit matrix
  $A \in \leanname{Transpose.Matrix}$ into $\lean{GridState}\;F\;d\;b \cong \Mat_{d\times d}(F^b)$
  ($b = \leanname{Transpose.entryBits}(d)$) via $0/1 \in F$ on each bit lane.
\end{enumerate}

\begin{proposition}[Equivariance of the native matrix-to-grid bridge]\label{prop:tm-grid-bridge}
\formal{DiagnosisAndBridgesConsequences.lean}{tmMatrix_gridLinear_bridges}
For every ring $F$, \lean{encodeTMMatrixToGrid} intertwines
\lean{stoneShuffleTMMatrix} with \lean{stoneShuffleGridLinear} (for $d = 2^m$),
\lean{stoneUnshuffleTMMatrix} with \lean{stoneUnshuffleGridLinear} (for $d = 2^m$, paired with its
linear-time $2$-work-tape machine \lean{stoneUnshuffleTM_twoTape_linear}),
\lean{rowShearTMMatrix} with \lean{rowShearLinear} (for $d \ge 1$), and
\leanname{Transpose.transpose} with \lean{gridSwapLinear}.
\end{proposition}
\begin{outline}
Verified directly by unfolding the $0/1$ grid embedding \lean{encodeTMMatrixToGrid} and checking
that cell/lane permutations commute with $\{0,1\} \hookrightarrow F$.
\end{outline}

\subsection{Formalized benchmarks from the Machine and Shared-Decoder Spines}\label{sec:benchmarks}

Before proving our new bounds, we formalize inside \texttt{BoundedArityDecoders} the exact
quantitative benchmarks from the Machine Spine that motivate reuniting $\Delta = O(1)$ with
decoder sharing. Doing so keeps the distinctions between word-step cones, bit-step cones, and
Paper~18's conditional adaptive interior-lane cones completely explicit.

\begin{proposition}[Bounded-cone reductions across trace models]\label{prop:diag-cones}
\formal{DiagnosisAndBarriersProofs.lean}{diagnosis_TCPAndAdaptiveConeReduction}
\formal{DiagnosisAndBridgesConsequences.lean}{tmWholeCellCutConeBound,wordStepTrace_TCPCellConeReduction,diagnosis_Paper18AdaptiveOriginalEntryPullback}
\formal{DiagnosisAndBarriers.lean}{Diagnosis_TCPAndAdaptiveConeReductionSpec}
\begin{enumerate}[label=(\alph*)]
\item \emph{$\TCP(k)$ walk-DAG and word-step cone compression}
  (\cite[Proposition~2.4 and Theorem~7.5]{BoundedStack2026}): for every $k \ge 1$ and $\eps > 0$,
  there exists a constant depth $D = D(k,\eps) \ge 1$ such that every $T$-step $k$-head unit-speed
  walk admits a cut step set $\Gamma \subseteq \{0,\dots,T-1\}$ of size $\abs{\Gamma} \le \eps T$
  reducing all $\Gamma$-free causal paths to $\le D$ vertices and all $\Gamma$-truncated causal
  cones to $\le (k+1)^D$ steps; on a word-step frozen trace (\lean{WordStepFrozenTrace}),
  $D = \lean{wordStepTCPDepth}(k,\eps)$ bounds every output cell's input-cell cone by $(k+1)^D$
  (\lean{wordStepTrace_TCPCellConeReduction}).
\item \emph{Bit-level \leanname{Transpose.TM} fibres and whole-cell cones}: for every $c_0 L$-step
  transposer $M$ ($1 \le M.k$) and $0 < \eps < 1$, there exist $D \ge 1$ and $D_0$ such that for all
  $d \ge \max(D_0, 2)$, a nonlinear fibre $C \subseteq \Sigma_d^{d\times d}$ of size
  $\abs{C} \ge 2^{(1-\eps)L}$ determines every output bit outside $\le \eps L$ broken bits from at
  most $(M.k+1)^D$ input cells; on any frozen trace/cut fibre of path length $\le D$, the $b$ bits
  of an output cell $(r,c)$ jointly read at most $b(M.k+1)^D$ input cells
  (\lean{tmWholeCellCutConeBound}).
\item \emph{Conditional adaptive interior-lane cut pullback}
  (\cite[Lemma~H.2, Theorem~H.5, and Corollary~H.6]{AdaptiveTransfer2026}): alongside the per-entry
  window-certificate Markov tail, for any adaptive transposer ($24 s \le b$) on a common
  simulator-trace fibre $S$, \emph{given} an original vertex cut $B$ whose $B$-free paths have at
  most $D$ vertices, $B$ bounds original-entry capacity with whole-entry cone $\le b(M.k+3)^D$ and
  transports to a simulator guard $\Gamma$ of size $\le G_s\abs{B}$ outside which every retained
  interior bit has simulator input-bit cone $\le 2(6s+1)(M.k+3)^D$
  (\lean{diagnosis_Paper18AdaptiveOriginalEntryPullback}).
\end{enumerate}
\end{proposition}
\begin{outline}
Part~(a) invokes \leanname{BoundedStackDepth.genie} (\cite{BoundedStack2026}) on the $k$-head walk
DAG and inducts over the $(k+1)$-ary dependency tree. Part~(b) pigeonholes the aligned trace/cut
advice on $c_0 L$-step runs and takes the $b$-fold union of bit-level cones. Part~(c) bridges to
\leanname{AdaptiveTransfer.TransferEntry.uniform_entry_success_tail},
\leanname{AdaptiveTransfer.TransferEntry.native_original_entry_capacity}, and
\leanname{AdaptiveTransfer.TransferCones.native_cone_transport_interior}.
\end{outline}

\begin{proposition}[Why cones without sharing stall: separated storage codes and barrel-shift networks]\label{prop:diag-high-rate}
\formal{DiagnosisAndBarriersProofs.lean}{diagnosis_AbstractSeparatedHighRateCode,diagnosis_ShiftAndRotationCodingGainBarrier}
\formal{DiagnosisAndBridgesConsequences.lean}{diagnosis2_HighRateTendstoOne,diagnosis3_MultiBitBlockEdgeShiftBarrier}
\formal{DiagnosisAndBarriers.lean}{Diagnosis_AbstractSeparatedHighRateCodeSpec,Diagnosis_ShiftAndRotationCodingGainBarrierSpec}
\begin{enumerate}[label=(\alph*)]
\item \emph{High-rate binary storage codes under paired-distance separation}
  (\cite[Theorem~7.1 and Corollary~8.2]{StorageCodes2026}; \cite[Definition~5.1, Theorem~5.5, and
  Theorem~6.1]{HighRate2026}): the Goren--Oort numerical parameters $(p, k) = (127, 91)$ exceed half
  rate both over $\F_{127^2}$ ($8100/16130 > 1/2$) and after $14$-bit binary blow-up ($> 0.5013$);
  moreover, given any finite field $K$ of odd characteristic ($(2 : K) \ne 0$) and odd-girth linear
  seed \lean{HasOddGirthLinearSeed} $K$ $D_{\mathrm{seed}}$ $\rho$ with $\frac{1}{2} < \rho < 1$,
  there exist $C_{\mathrm{code}} > 0$, $0 < \alpha < \frac{1}{2}$, and $\Delta_0$ such that every
  degree cap $\Delta \ge \Delta_0$ and paired-distance girth threshold $G$ admit a nonempty
  $(\Delta, G)$-separated binary storage code $C \subseteq \{0,1\}^n$ of rate
  $\frac{\log_2\abs{C}}{n} \ge 1 - C_{\mathrm{code}}\Delta^{-\alpha} \to 1$ as $\Delta \to \infty$
  (\lean{diagnosis2_HighRateTendstoOne}).
\item \emph{Barrel-shift complete-exchange coding-gain barrier}
  (\cite[Propositions~11.2 and~11.3]{FirstPaper2026}): any size-$s$ fan-in-$2$ acyclic Boolean
  barrel-shift circuit on $n \ge N_0$ bits (\lean{HasBarrelShiftCircuit}) embeds into a max-degree-$4$
  graph on $V \le C_{\mathrm{net}}(s + n)$ vertices carrying an $n$-round causal complete-exchange
  code (\lean{HasCausalCompleteExchangeCode}) whose coding-to-routing gain against any multicommodity
  routing rate $\lambda_{\mathrm{route}} > 0$ is at least
  $c_{\mathrm{gain}}\,n\log_2 n / (s + n)$ (and $\ge c_{\mathrm{gain}}\,n\log_2 n / (w s_{\mathrm{blk}} + n)$
  on $s_{\mathrm{blk}}$ $w$-bit block edges, \lean{diagnosis3_MultiBitBlockEdgeShiftBarrier}), while
  every $n \ge 2$ admits an explicit barrel shifter of size $s \le C_{\mathrm{sh}} n\log_2 n$.
\end{enumerate}
\end{proposition}
\begin{outline}
Part~(a) verifies the $(127, 91)$ arithmetic inequalities, amplifies the odd-characteristic linear
seed by recursive odd-majority tensor products on syndrome fibres, and takes $\Delta \to \infty$.
Part~(b) converts each fan-in-$2$ gate into a bounded-degree DAG gadget serving all $n$ cyclic shift
demands in $n$ rounds, bounds pairwise source-to-sink walk lengths via the $n!$ shift permutations
($\sum_{i,j}\mathrm{distBound}(i,j) \ge c_{\mathrm{dist}} n^2\log_2 n$), and builds the
$O(n\log n)$ barrel shifter from $\clog{n}$ power-of-two multiplexor layers.
\end{outline}
